% % --- 题目 ---
% \problem[P227-8 (98-3)]{设 $\displaystyle X_1, X_2, X_3, X_4$ 为来自正态总体 $\displaystyle N(0, 2^2)$ 的简单随机样本;
% \[ X = a(X_1 - 2X_2)^2 + b(3X_3 - 4X_4)^2, \]
% 其中 $\displaystyle a, b \neq 0$. 则当 $\displaystyle a = \underline{\hspace{4em}}, b = \underline{\hspace{4em}}$ 时, 统计量 $\displaystyle X$ 服从 $\displaystyle \chi^2$ 分布, 其自由度为 \underline{\hspace{4em}}.}
% \ansat{370；【真题精选】 - 考点一：抽样分布 - 8}

% % --- 题目 ---
% \problem[P200-2 (16-1)]{设随机变量 $\displaystyle X \sim N(\mu, \sigma^2) \ (\sigma > 0)$, 记 $\displaystyle p = P\{X \le \mu + \sigma^2\}$, 则~(~\quad~)}
% \begin{tasks}(2)
%   \task $\displaystyle p$ 随着 $\displaystyle \mu$ 的增加而增加.
%   \task $\displaystyle p$ 随着 $\displaystyle \sigma$ 的增加而增加.
%   \task $\displaystyle p$ 随着 $\displaystyle \mu$ 的增加而减少.
%   \task $\displaystyle p$ 随着 $\displaystyle \sigma$ 的增加而减少.
% \end{tasks}
% \ansat{355；【十年真题】 - 考点一：随机变量的分布 - 2}

% % --- 题目 ---
% \problem[P203-4 (02-1)]{设 $\displaystyle X_1$ 和 $\displaystyle X_2$ 是相互独立的连续型随机变量, 它们的密度函数分别为 $\displaystyle f_1(x)$ 和 $\displaystyle f_2(x)$, 分布函数分别为 $\displaystyle F_1(x)$ 和 $\displaystyle F_2(x)$, 则~(~\quad~)}
% \begin{tasks}(1)
%   \task $\displaystyle f_1(x) + f_2(x)$ 必为某一随机变量的概率密度.
%   \task $\displaystyle f_1(x) f_2(x)$ 必为某一随机变量的概率密度.
%   \task $\displaystyle F_1(x) + F_2(x)$ 必为某一随机变量的分布函数.
%   \task $\displaystyle F_1(x) F_2(x)$ 必为某一随机变量的分布函数.
% \end{tasks}
% \ansat{356；【真题精选】 - 考点一：随机变量的分布 - 4}

% % --- 题目 ---
% \problem[P219-3 (13-3)]{设随机变量 $\displaystyle X$ 服从标准正态分布 $\displaystyle N(0, 1)$, 则 $\displaystyle E(X\mathrm{e}^{2X}) = \underline{\hspace{4em}}.$}
% \ansat{366；【真题精选】 - 考点一：随机变量的数学期望与方差 - 3}
% \vspace{6em}

% % --- 题目 ---
% \problem[P227-2 (11-3)]{设总体 $\displaystyle X$ 服从参数为 $\displaystyle \lambda (\lambda > 0)$ 的泊松分布, $\displaystyle X_1, X_2, \dots, X_n (n \ge 2)$ 为来自总体 $\displaystyle X$ 的简单随机样本, 则对应的统计量 $\displaystyle T_1 = \frac{1}{n} \sum_{i=1}^n X_i, T_2 = \frac{1}{n-1} \sum_{i=1}^{n-1} X_i + \frac{1}{n} X_n$, 有 ( \quad )
% \begin{tasks}(2)
%   \task $\displaystyle ET_1 > ET_2, DT_1 > DT_2$.
%   \task $\displaystyle ET_1 > ET_2, DT_1 < DT_2$.
%   \task $\displaystyle ET_1 < ET_2, DT_1 > DT_2$.
%   \task $\displaystyle ET_1 < ET_2, DT_1 < DT_2$.
% \end{tasks}}
% \ansat{370；【真题精选】 - 考点二：统计量的数字特征 - 2}

% % --- 题目 ---
% \problem[P224-1 (23-3-仅数学三)]{设 $\displaystyle X_1, X_2$ 为来自总体 $\displaystyle N(\mu, \sigma^2)$ 的简单随机样本, 其中 $\displaystyle \sigma (\sigma > 0)$ 是未知参数. 记 $\displaystyle \hat{\sigma} = a|X_1 - X_2|$, 若 $\displaystyle E(\hat{\sigma}) = \sigma$, 则 $\displaystyle a = ( \quad )$
% \begin{tasks}(4)
%   \task $\displaystyle \frac{\sqrt{\pi}}{2}$.
%   \task $\displaystyle \frac{\sqrt{2\pi}}{2}$.
%   \task $\displaystyle \sqrt{\pi}$.
%   \task $\displaystyle \sqrt{2\pi}$.
% \end{tasks}}
% \ansat{370；【十年真题】 - 考点二：统计量的数字特征 - 1}

% --- 题目 ---
\problem[P215-13 (23-1)]{设二维随机变量 $\displaystyle (X, Y)$ 的概率密度为
\[ f(x, y) = \begin{cases} \displaystyle \frac{2}{\pi}(x^2 + y^2), & \displaystyle x^2 + y^2 \le 1, \\ 0, & \text{其他}. \end{cases} \]
\begin{enumerate}
\item[(1)] 求 $\displaystyle X$ 与 $\displaystyle Y$ 的协方差;
\item[(2)] $\displaystyle X$ 与 $\displaystyle Y$ 是否相互独立?
\item[(3)] 求 $\displaystyle Z = X^2 + Y^2$ 的概率密度.
\end{enumerate}}
\begin{note}
  主要错的是 (2)
\end{note}
\ansat{365；【十年真题】 - 考点二：随机变量的协方差与相关系数 - 13}

% --- 题目 ---
\problem[P229-2 (22-1,3)]{设 $\displaystyle X_1, X_2, \dots, X_n$ 为来自均值为 $\displaystyle \theta$ 的指数分布总体的简单随机样本, $\displaystyle Y_1, Y_2, \dots, Y_m$ 为来自均值为 $\displaystyle 2\theta$ 的指数分布总体的简单随机样本, 且两样本相互独立, 其中 $\displaystyle \theta (\theta > 0)$ 为未知参数. 利用样本 $\displaystyle X_1, X_2, \dots, X_n, Y_1, Y_2, \dots, Y_m$, 求 $\displaystyle \theta$ 的最大似然估计量 $\displaystyle \hat{\theta}$, 并求 $\displaystyle D(\hat{\theta})$.
}
\ansat{371；【十年真题】 - 考点一：矩估计与最大似然估计 - 2}

% --- 题目 ---
\problem[P219-12 (15-1,3)]{设随机变量 $\displaystyle X$ 的概率密度为
\[ f(x) = \begin{cases} 2^{-x} \ln 2, & x > 0, \\ 0, & x \le 0. \end{cases} \]
对 $\displaystyle X$ 进行独立重复的观测, 直到第 2 个大于 3 的观测值出现时停止, 记 $\displaystyle Y$ 为观测次数.
\begin{enumerate}
\item[(1)] 求 $\displaystyle Y$ 的概率分布;
\item[(2)] 求 $\displaystyle EY$.
\end{enumerate}}
\begin{note}
  主要错的是 (2)
\end{note}
\ansat{366；【真题精选】 - 考点一：随机变量的数学期望与方差 - 12}

% --- 题目 ---
\problem[P198-11 (97-1)]{袋中有 $\displaystyle 50$ 个乒乓球, 其中 $\displaystyle 20$ 个是黄球, $\displaystyle 30$ 个是白球, 今有两人依次随机地从袋中各取一球, 取后不放回, 则第二个人取得黄球的概率是 \underline{\hspace{4em}}.}
\ansat{354；【真题精选】 - 考点一：概率的五大公式 - 11}

% % --- 题目 ---
% \problem[P220-14 (03-1)]{已知甲、乙两箱中装有同种产品, 其中甲箱中装有 3 件合格品和 3 件次品, 乙箱中仅装有 3 件合格品. 从甲箱中任取 3 件产品放入乙箱后, 求:
% \begin{enumerate}
% \item[(1)] 乙箱中次品件数 $\displaystyle X$ 的数学期望;
% \item[(2)] 从乙箱中任取一件产品是次品的概率.
% \end{enumerate}}
% \ansat{366；【真题精选】 - 考点一：随机变量的数学期望与方差 - 14}

% --- 题目 ---
\problem[P215-10 (23-3)]{设随机变量 $\displaystyle X$ 与 $\displaystyle Y$ 相互独立, 且 $\displaystyle X \sim B(1, p)$, $\displaystyle Y \sim B(2, p), p \in (0, 1)$, 则 $\displaystyle X+Y$ 与 $\displaystyle X-Y$ 的相关系数为 \underline{\hspace{4em}}.}
\ansat{364；【十年真题】 - 考点二：随机变量的协方差与相关系数 - 10}

% --- 题目 ---
\problem[P215-8 (20-3)]{设随机变量 $\displaystyle (X, Y)$ 服从二维正态分布 $\displaystyle N\left(0, 0; 1, 4; -\frac{1}{2}\right)$, 则下列随机变量中服从标准正态分布且与 $\displaystyle X$ 独立的是~(~\quad~)}
\begin{tasks}(4)
  \task $\displaystyle \frac{\sqrt{5}}{5}(X+Y).$
  \task $\displaystyle \frac{\sqrt{5}}{5}(X-Y).$
  \task $\displaystyle \frac{\sqrt{3}}{3}(X+Y).$
  \task $\displaystyle \frac{\sqrt{3}}{3}(X-Y).$
\end{tasks}
\ansat{364；【十年真题】 - 考点二：随机变量的协方差与相关系数 - 8}

% % --- 题目 ---
% \problem[P216-15 (19-1,3)]{设随机变量 $\displaystyle X$ 与 $\displaystyle Y$ 相互独立, $\displaystyle X$ 服从参数为 1 的指数分布, $\displaystyle Y$ 的概率分布为
% \[ P\{Y = -1\} = p, \quad P\{Y = 1\} = 1-p \quad (0 < p < 1). \]
% 令 $\displaystyle Z = XY$.
% \begin{enumerate}
% \item[(1)] 求 $\displaystyle Z$ 的概率密度;
% \item[(2)] $\displaystyle p$ 为何值时, $\displaystyle X$ 与 $\displaystyle Z$ 不相关?
% \item[(3)] $\displaystyle X$ 与 $\displaystyle Z$ 是否相互独立?
% \end{enumerate}}
% \ansat{365；【十年真题】 - 考点二：随机变量的协方差与相关系数 - 15}

% % --- 题目 ---
% \problem[P211-9 (95-3)]{已知随机变量 $\displaystyle X$ 和 $\displaystyle Y$ 的联合概率密度为
% \[ \varphi(x, y) = \begin{cases} 4xy, & 0 \le x \le 1, 0 \le y \le 1, \\ 0, & \text{其他}, \end{cases} \]
% 求 $\displaystyle X$ 和 $\displaystyle Y$ 联合分布函数 $\displaystyle F(x, y)$.}
% \ansat{360；【真题精选】 - 考点一：二维随机变量的分布 - 9}

% % --- 题目 ---
% \problem[P200-3 (24-1,3)]{设随机试验每次成功的概率为 $\displaystyle p$, 现进行 3 次独立重复试验, 在至少成功 1 次的条件下 3 次试验全部成功的概率为 $\frac{4}{13}$, 则 $\displaystyle p = \underline{\hspace{4em}}. $}
% \ansat{355；【十年真题】 - 考点一：随机变量的分布 - 3}

% % --- 题目 ---
% \problem[P219-2 (09-1)]{设随机变量 $\displaystyle X$ 的分布函数为 $\displaystyle F(x) = 0.3\Phi(x) + 0.7\Phi\left(\frac{x-1}{2}\right)$, 其中 $\displaystyle \Phi(x)$ 为标准正态分布函数, 则 $\displaystyle EX $= ( \quad )}
% \begin{tasks}(4)
%   \task 0.
%   \task 0.3.
%   \task 0.7.
%   \task 1.
% \end{tasks}
% \ansat{366；【真题精选】 - 考点一：随机变量的数学期望与方差 - 2}
% \vspace{8em}

% % --- 题目 ---
% \problem[P211-7 (01-1)]{设某班车起点站上车人数 $\displaystyle X$ 服从参数为 $\displaystyle \lambda \ (\lambda > 0)$ 的泊松分布, 每位乘客在中途下车的概率为 $\displaystyle p \ (0 < p < 1)$, 且中途下车与否相互独立. 以 $\displaystyle Y$ 表示在中途下车的人数, 求:
% \begin{enumerate}
% \item[(1)] 在发车时有 $\displaystyle n$ 个乘客的条件下, 中途有 $\displaystyle m$ 人下车的概率.
% \item[(2)] 二维随机变量 $\displaystyle (X, Y)$ 的概率分布.
% \end{enumerate}}
% \ansat{360；【真题精选】 - 考点一：二维随机变量的分布 - 7}

% --- 题目 ---
\problem[P222-1 (22-1)]{设随机变量 $\displaystyle X_1, X_2, \dots, X_n$ 独立同分布. 且 $\displaystyle X_1$ 的 $\displaystyle 4$ 阶矩存在. 记 $\displaystyle \mu_k = E(X_1^k) \ (k=1,2,3,4)$, 则根据切比雪夫不等式, 对任意 $\displaystyle \varepsilon > 0$, 有
\[ P\left\{\left|\frac{1}{n}\sum_{i=1}^n X_i^2 - \mu_2\right| \ge \varepsilon\right\} \le (\quad) \]
}
\begin{tasks}(4)
  \task $\displaystyle \frac{\mu_4 - \mu_2^2}{n\varepsilon^2}$.
  \task $\displaystyle \frac{\mu_4 - \mu_2^2}{\sqrt{n}\varepsilon^2}$.
  \task $\displaystyle \frac{\mu_2 - \mu_1^2}{n\varepsilon^2}$.
  \task $\displaystyle \frac{\mu_2 - \mu_1^2}{\sqrt{n}\varepsilon^2}$.
\end{tasks}
\ansat{368；【十年真题】 - 考点：大数定律、中心极限定理与切比雪夫不等式 - 1}

% % --- 题目 ---
% \problem[P216-例2]{设随机变量 $\displaystyle X \sim N(0, 1), Y \sim N(1, 4)$, 且 $\displaystyle U = aX + Y, V = aX - Y (a > 0)$. 若 $\displaystyle U, V$ 不相关, 则 $\displaystyle a = \underline{\hspace{4em}}.$}
% \ansat{219；【方法探究】 - 考点二：随机变量的协方差与相关系数 - 例2}

% % --- 题目 ---
% \problem[P204-11 (89-4)]{设随机变量 $\displaystyle X$ 的分布函数为
% \[ F(x) = \begin{cases} 0, & \displaystyle x < 0, \\ \displaystyle A \sin x, & \displaystyle 0 \le x \le \frac{\pi}{2}, \\ 1, & \displaystyle x > \frac{\pi}{2}, \end{cases} \]
% 则 $\displaystyle P\left\{|X| < \frac{\pi}{6}\right\} = \underline{\hspace{4em}}. $}
% \ansat{356；【真题精选】 - 考点一：随机变量的分布 - 11}

% % --- 题目 ---
% \problem[P204-2 (03-3)]{设随机变量 $\displaystyle X$ 的概率密度为
% $$ f(x) = \begin{cases} \displaystyle \frac{1}{3\sqrt[3]{x^2}}, & x \in [1, 8], \\ 0, & \text{其他}, \end{cases} $$
% $\displaystyle F(x)$ 是 $\displaystyle X$ 的分布函数, 则随机变量 $\displaystyle Y=F(X)$ 的分布函数为 \underline{\hspace{4em}}.}
% \ansat{357；【真题精选】 - 考点二：随机变量的函数的分布 - 2}

% --- 题目 ---
\problem[P196-7 (25-1)]{设 $\displaystyle A, B$ 为两个随机事件, 且 $\displaystyle A$ 与 $\displaystyle B$ 相互独立. 已知 $\displaystyle P(A)=2P(B)$, $\displaystyle P(A \cup B)=\frac{5}{8}$, 则在事件 $\displaystyle A, B$ 至少有一个发生的条件下, $\displaystyle A, B$ 中恰有一个发生的概率为 \underline{\hspace{4em}}.}
\ansat{353；【十年真题】 - 考点一：概率的五大公式 - 7}

% --- 题目 ---
\problem[P202-例3 (90-1)]{已知随机变量 $\displaystyle X$ 的概率密度函数 $\displaystyle f(x) = \frac{1}{2}\mathrm{e}^{-|x|}, -\infty < x < +\infty$, 则 $\displaystyle X$ 的分布函数 $\displaystyle F(x) = \underline{\hspace{4em}}. $}
\ansat{202；【方法探究】 - 考点一：随机变量的分布 - 例3}
\vspace{2em}

% % --- 题目 ---
% \problem[P216-16 (18-1,3)]{设随机变量 $\displaystyle X$ 与 $\displaystyle Y$ 相互独立, $\displaystyle X$ 的概率分布为
% \[ P\{X=1\} = P\{X=-1\} = \frac{1}{2}, \]
% $\displaystyle Y$ 服从参数为 $\displaystyle \lambda$ 的泊松分布. 令 $\displaystyle Z=XY$.
% \begin{enumerate}
% \item[(1)] 求 $\displaystyle \mathrm{Cov}(X, Z)$;
% \item[(2)] 求 $\displaystyle Z$ 的概率分布.
% \end{enumerate}}
% \ansat{365；【十年真题】 - 考点二：随机变量的协方差与相关系数 - 16}

% --- 题目 ---
\problem[P210-4 (10-1,3)]{设二维随机变量 $\displaystyle (X, Y)$ 的概率密度为 $\displaystyle f(x, y) = A\mathrm{e}^{-2x^2 + 2xy - y^2}, -\infty < x < +\infty, -\infty < y < +\infty$, 求常数 $\displaystyle A$ 及条件概率密度 $\displaystyle f_{Y|X}(y|x)$.}
\ansat{360；【真题精选】 - 考点一：二维随机变量的分布 - 4}

% % --- 题目 ---
% \problem[P214-1 (24-3)]{设随机变量 $\displaystyle X$ 的概率密度为
% \[ f(x) = \begin{cases} 6x(1-x), & 0 < x < 1, \\ 0, & \text{其他}, \end{cases} \]
% 则 $\displaystyle X$ 的三阶中心矩 $\displaystyle E[(X-EX)^3] = (\quad)$
% }
% \begin{tasks}(4)
%   \task $\displaystyle -\frac{1}{32}.$
%   \task $\displaystyle 0.$
%   \task $\displaystyle \frac{1}{16}.$
%   \task $\displaystyle \frac{1}{2}.$
% \end{tasks}
% \ansat{362；【十年真题】 - 考点一：随机变量的数学期望与方差 - 1}

% % --- 题目 ---
% \problem[P227-7 (99-3)]{在天平上重复称量一重为 $\displaystyle a$ 的物品, 假设各次称量结果相互独立且都服从正态分布 $\displaystyle N(a, 0.2^2)$. 若以 $\displaystyle \bar{X}_n$ 表示 $\displaystyle n$ 次称量结果的算术平均值, 则为使 $\displaystyle P\{|\bar{X}_n - a| < 0.1\} \ge 0.95$, $\displaystyle n$ 的最小值应不小于自然数 \underline{\hspace{4em}}.
% (注: 标准正态分布函数值 $\displaystyle \Phi(1.96) = 0.975$)}
% \ansat{370；【真题精选】 - 考点一：抽样分布 - 7}

% % --- 题目 ---
% \problem[P214-10 (21-1,3)]{在区间 $\displaystyle (0, 2)$ 上随机取一点, 将该区间分成两段, 较短一段的长度记为 $\displaystyle X$, 较长一段的长度记为 $\displaystyle Y$. 令 $\displaystyle Z = \frac{Y}{X}$.
% \begin{enumerate}
% \item[(1)] 求 $\displaystyle X$ 的概率密度;
% \item[(2)] 求 $\displaystyle Z$ 的概率密度;
% \item[(3)] 求 $\displaystyle E\left(\frac{X}{Y}\right)$.
% \end{enumerate}}
% \begin{note}
%   主要错的是 (3)
% \end{note}
% \ansat{363；【十年真题】 - 考点一：随机变量的数学期望与方差 - 10}

% --- 题目 ---
\problem[P206-1 (24-1,3)]{设随机变量 $\displaystyle X, Y$ 相互独立, 且均服从参数为 $\displaystyle \lambda$ 的指数分布. 令 $\displaystyle Z = |X - Y|$, 则下列随机变量中与 $\displaystyle Z$ 同分布的是~(~\quad~)}
\begin{tasks}(4)
  \task $\displaystyle X + Y.$
  \task $\displaystyle \frac{X+Y}{2}.$
  \task $\displaystyle 2X.$
  \task $\displaystyle X.$
\end{tasks}
\ansat{358；【十年真题】 - 考点二：两个随机变量的函数的分布 - 1}

% --- 题目 ---
\problem[P215-12 (20-1)]{设 $\displaystyle X$ 服从区间 $\displaystyle \left(-\frac{\pi}{2}, \frac{\pi}{2}\right)$ 上的均匀分布, $\displaystyle Y = \sin X$, 则 $\displaystyle \mathrm{Cov}(X, Y) = \underline{\hspace{4em}}. $}
\ansat{365；【十年真题】 - 考点二：随机变量的协方差与相关系数 - 12}

% --- 题目 ---
\problem[P224-2 (25-3-仅数学三)]{设总体 $\displaystyle X$ 的分布函数为 $\displaystyle F(x)$, $\displaystyle X_1, X_2, \dots, X_n$ 为来自总体 $\displaystyle X$ 的简单随机样本, 样本的经验分布函数为 $\displaystyle F_n(x)$. 对于给定的 $\displaystyle x$ ($\displaystyle 0 < F(x) < 1$), $\displaystyle D[F_n(x)] = ( \quad )$
\begin{tasks}(2)
  \task $\displaystyle F(x)[1-F(x)]$.
  \task $\displaystyle [F(x)]^2$.
  \task $\displaystyle \frac{1}{n} F(x)[1-F(x)]$.
  \task $\displaystyle \frac{1}{n} [F(x)]^2$.
\end{tasks}}
\ansat{369；【十年真题】 - 考点一：抽样分布 - 2}

% --- 题目 ---
\problem[P199-14 (89-1)]{甲、乙两人独立地对同一目标射击一次, 其命中率分别为 $\displaystyle 0.6$ 和 $\displaystyle 0.5$. 现已知目标被命中, 则它是甲射中的概率为 \underline{\hspace{4em}}.}
\ansat{355；【真题精选】 - 考点一：概率的五大公式 - 14}
\vspace{6em}

% % --- 题目 ---
% \problem[P229-4 (19-1,3)]{设总体 $\displaystyle X$ 的概率密度为
% \[ f(x; \sigma^2) = \begin{cases} \displaystyle \frac{A}{\sigma} \mathrm{e}^{-\frac{(x-\mu)^2}{2\sigma^2}}, & x \ge \mu, \\ 0, & x < \mu, \end{cases} \]
% 其中 $\displaystyle \mu$ 是已知参数, $\displaystyle \sigma > 0$ 是未知参数, $\displaystyle A$ 是常数, $\displaystyle X_1, X_2, \dots, X_n$ 是来自总体 $\displaystyle X$ 的简单随机样本.
% \begin{enumerate}
% \item[(1)] 求 $\displaystyle A$;
% \item[(2)] 求 $\displaystyle \sigma^2$ 的最大似然估计量.
% \end{enumerate}}
% \ansat{372；【十年真题】 - 考点一：矩估计与最大似然估计 - 4}

% --- 题目 ---
\problem[P211-5 (09-1,3)]{袋中有 $\displaystyle 1$ 个红球, $\displaystyle 2$ 个黑球与 $\displaystyle 3$ 个白球, 现有放回地从袋中取两次, 每次取一个球. 以 $\displaystyle X, Y, Z$ 分别表示两次取球所取得的红球、黑球与白球的个数.
\begin{enumerate}
\item[(1)] 求 $\displaystyle P\{X=1|Z=0\}$;
\item[(2)] 求二维随机变量 $\displaystyle (X,Y)$ 的概率分布.
\end{enumerate}}
\ansat{360；【真题精选】 - 考点一：二维随机变量的分布 - 5}

% --- 题目 ---
\problem[P227-5 (08-1,3)]{设 $\displaystyle X_1, X_2, \dots, X_n$ 是总体 $\displaystyle N(\mu, \sigma^2)$ 的简单随机样本. 记 $\displaystyle \bar{X} = \frac{1}{n} \sum_{i=1}^n X_i, S^2 = \frac{1}{n-1} \sum_{i=1}^n (X_i - \bar{X})^2, \ T = \bar{X}^2 - \frac{1}{n} S^2$.
\begin{enumerate}
\item[(1)] 证明 $\displaystyle ET = \mu^2$;
\item[(2)] 当 $\displaystyle \mu=0, \sigma=1$ 时, 求 $\displaystyle DT$.
\end{enumerate}}
\begin{note}
  主要错的是 (2)
\end{note}
\ansat{371；【真题精选】 - 考点二：统计量的数字特征 - 5}

% --- 题目 ---
\problem[P215-5 (22-1)]{设随机变量 $\displaystyle X \sim N(0, 1)$, 在 $\displaystyle X=x$ 的条件下随机变量 $\displaystyle Y \sim N(x, 1)$, 则 $\displaystyle X$ 与 $\displaystyle Y$ 的相关系数为~(~\quad~)}
\begin{tasks}(4)
  \task $\displaystyle \frac{1}{4}.$
  \task $\displaystyle \frac{1}{2}.$
  \task $\displaystyle \frac{\sqrt{3}}{3}.$
  \task $\displaystyle \frac{\sqrt{2}}{2}.$
\end{tasks}
\ansat{364；【十年真题】 - 考点二：随机变量的协方差与相关系数 - 5}

% % --- 题目 ---
% \problem[P223-例]{设随机变量 $\displaystyle X$ 服从二项分布 $\displaystyle B(100, 0.2)$, 则利用中心极限定理可得 $\displaystyle P\{X \ge 12\}$ 的近似值为 ( \quad )}
% \begin{tasks}(4)
%   \task $\displaystyle \Phi(2)$.
%   \task $\displaystyle 1-\Phi(2)$.
%   \task $\displaystyle \Phi(0.5)$.
%   \task $\displaystyle 1-\Phi(0.5)$.
% \end{tasks}
% \ansat{223；【方法探究】 - 考点：大数定律、中心极限定理与切比雪夫不等式 - 例}

% --- 题目 ---
\problem[P229-3 (20-1,3)]{设某种元件的使用寿命 $\displaystyle T$ 的分布函数为
\[ F(t) = \begin{cases} 1 - \mathrm{e}^{-(\frac{t}{\theta})^m}, & t \ge 0, \\ 0, & \text{其他}, \end{cases} \]
其中 $\displaystyle \theta, m$ 为参数且大于零.
\begin{enumerate}
\item[(1)] 求概率 $\displaystyle P\{T>t\}$ 与 $\displaystyle P\{T>t+s | T>s\}$, 其中 $\displaystyle s>0, t>0$;
\item[(2)] 任取 $\displaystyle n$ 个这种元件做寿命试验, 测得它们的寿命分别为 $\displaystyle t_1, t_2, \dots, t_n$. 若 $\displaystyle m$ 已知, 求 $\displaystyle \theta$ 的最大似然估计值 $\displaystyle \hat{\theta}$.
\end{enumerate}}
\begin{note}
  主要错的是 (2)
\end{note}
\ansat{371；【十年真题】 - 考点一：矩估计与最大似然估计 - 3}

% --- 题目 ---
\problem[P214-8 (25-1,3)]{投保人的损失事件发生时, 保险公司的赔付额 $\displaystyle Y$ 与投保人的损失额 $\displaystyle X$ 的关系为
\[ Y = \begin{cases} 0, & X \le 100, \\ X - 100, & X > 100. \end{cases} \]
设损失事件发生时, 投保人的损失额 $\displaystyle X$ 的概率密度为
\[ f(x) = \begin{cases} \displaystyle \frac{2 \times 100^2}{(100+x)^3}, & x > 0, \\ 0, & x \le 0. \end{cases} \]
\begin{enumerate}
\item[(1)] 求 $\displaystyle P\{Y > 0\}$ 及 $\displaystyle EY$;
\item[(2)] 这种损失事件在一年内发生的次数记为 $\displaystyle N$, 保险公司在一年内就这种损失事件产生的理赔次数记为 $\displaystyle M$. 假设 $\displaystyle N$ 服从参数为 8 的泊松分布, 在 $\displaystyle N=n \ (n \ge 1)$ 的条件下, $\displaystyle M$ 服从二项分布 $\displaystyle B(n, p)$, 其中 $\displaystyle p=P\{Y > 0\}$. 求 $\displaystyle M$ 的概率分布.
\end{enumerate}}
\begin{note}
  主要错的是 (2)
\end{note}
\ansat{363；【十年真题】 - 考点一：随机变量的数学期望与方差 - 8}

% % --- 题目 ---
% \problem[P204-7 (93-3)]{设随机变量 $\displaystyle X$ 的概率密度为 $\displaystyle \varphi(x)$, 且 $\displaystyle \varphi(-x) = \varphi(x)$. $\displaystyle F(x)$ 为 $\displaystyle X$ 的分布函数, 则对任意实数 $\displaystyle a$, 有~(~\quad~)}
% \begin{tasks}(2)
%   \task $\displaystyle F(-a) = 1 - \int_{0}^{a} \varphi(x) \ \mathrm{d}x.$
%   \task $\displaystyle F(-a) = \frac{1}{2} - \int_{0}^{a} \varphi(x) \ \mathrm{d}x.$
%   \task $\displaystyle F(-a) = F(a).$
%   \task $\displaystyle F(-a) = 2F(a) - 1.$
% \end{tasks}
% \ansat{356；【真题精选】 - 考点一：随机变量的分布 - 7}

% % --- 题目 ---
% \problem[P224-4 (18-3)]{设 $\displaystyle X_1, X_2, \dots, X_n (n \ge 2)$ 为来自总体 $\displaystyle N(\mu, \sigma^2)$ ($\displaystyle \sigma > 0$) 的简单随机样本. 令 $\displaystyle \bar{X} = \frac{1}{n} \sum_{i=1}^n X_i, \ S = \sqrt{\frac{1}{n-1} \sum_{i=1}^n (X_i - \bar{X})^2}, S^* = \sqrt{\frac{1}{n} \sum_{i=1}^n (X_i - \mu)^2},$
% 则 ( \quad )
% \begin{tasks}(2)
%   \task $\displaystyle \frac{\sqrt{n}(\bar{X} - \mu)}{S} \sim t(n)$.
%   \task $\displaystyle \frac{\sqrt{n}(\bar{X} - \mu)}{S} \sim t(n-1)$.
%   \task $\displaystyle \frac{\sqrt{n}(\bar{X} - \mu)}{S^*} \sim t(n)$.
%   \task $\displaystyle \frac{\sqrt{n}(\bar{X} - \mu)}{S^*} \sim t(n-1)$.
% \end{tasks}}
% \ansat{369；【十年真题】 - 考点一：抽样分布 - 4}
% \vspace{6em}

% % --- 题目 ---
% \problem[P206-3 (20-1)]{设随机变量 $\displaystyle X_1, X_2, X_3$ 相互独立, 其中 $\displaystyle X_1$ 与 $\displaystyle X_2$ 均服从标准正态分布, $\displaystyle X_3$ 的概率分布为 $\displaystyle P\{X_3=0\} = P\{X_3=1\} = \frac{1}{2}$. $\displaystyle Y = X_3 X_1 + (1-X_3)X_2$.
% \begin{enumerate}
% \item[(1)] 求二维随机变量 $\displaystyle (X_1, Y)$ 的分布函数, 结果用标准正态分布函数 $\displaystyle \Phi(x)$ 表示;
% \item[(2)] 证明随机变量 $\displaystyle Y$ 服从标准正态分布.
% \end{enumerate}}
% \begin{note}
%   主要错的是 (2)
% \end{note}
% \ansat{358；【十年真题】 - 考点二：两个随机变量的函数的分布 - 3}

% % --- 题目 ---
% \problem[P215-1 (25-1)]{设二维随机变量 $\displaystyle (X, Y)$ 服从正态分布 $\displaystyle N(0, 0; 1, 1; \rho)$, 其中 $\displaystyle \rho \in (-1, 1)$. 若 $\displaystyle a, b$ 为满足 $\displaystyle a^2 + b^2 = 1$ 的任意实数, 则 $\displaystyle D(aX + bY)$ 的最大值为~(~\quad~)}
% \begin{tasks}(4)
%   \task $\displaystyle 1.$
%   \task $\displaystyle 2.$
%   \task $\displaystyle 1 + |\rho|.$
%   \task $\displaystyle 1 + \rho^2.$
% \end{tasks}
% \ansat{363；【十年真题】 - 考点二：随机变量的协方差与相关系数 - 1}

% % --- 题目 ---
% \problem[P206-2 (23-1-局部)]{设二维随机变量 $\displaystyle (X, Y)$ 的概率密度为
% \[ f(x, y) = \begin{cases} \displaystyle \frac{2}{\pi}(x^2 + y^2), & \displaystyle x^2 + y^2 \le 1, \\ 0, & \text{其他}. \end{cases} \]
% \begin{enumerate}
% \item[(1)] $\displaystyle X$ 与 $\displaystyle Y$ 是否相互独立?
% \item[(2)] 求 $\displaystyle Z = X^2 + Y^2$ 的概率密度.
% \end{enumerate}}
% \ansat{358；【十年真题】 - 考点二：两个随机变量的函数的分布 - 2}

% % --- 题目 ---
% \problem[P198-1 (15-1,3)]{若 $\displaystyle A, B$ 为任意两个随机事件, 则~(~\quad~)}
% \begin{tasks}(2)
%   \task $\displaystyle P(AB) \le P(A)P(B).$
%   \task $\displaystyle P(AB) \ge P(A)P(B).$
%   \task $\displaystyle P(AB) \le \frac{P(A)+P(B)}{2}.$
%   \task $\displaystyle P(AB) \ge \frac{P(A)+P(B)}{2}.$
% \end{tasks}
% \ansat{354；【真题精选】 - 考点一：概率的五大公式 - 1}
% \vspace{4em}

% % --- 题目 ---
% \problem[P229-5 (18-1,3)]{设总体 $\displaystyle X$ 的概率密度为
% \[ f(x; \sigma) = \frac{1}{2\sigma} \mathrm{e}^{-\frac{|x|}{\sigma}}, \quad -\infty < x < +\infty, \]
% 其中 $\displaystyle \sigma \in (0, +\infty)$ 为未知参数, $\displaystyle X_1, X_2, \dots, X_n$ 为来自总体 $\displaystyle X$ 的简单随机样本. 记 $\displaystyle \sigma$ 的最大似然估计量为 $\displaystyle \hat{\sigma}$.
% \begin{enumerate}
% \item[(1)] 求 $\displaystyle \hat{\sigma}$;
% \item[(2)] 求 $\displaystyle E\hat{\sigma}$ 和 $\displaystyle D\hat{\sigma}$.
% \end{enumerate}}
% \ansat{372；【十年真题】 - 考点一：矩估计与最大似然估计 - 3}

% % --- 题目 ---
% \problem[P198-例(1)]{袋中装有 $\displaystyle 2$ 个红球, $\displaystyle 3$ 个黄球, $\displaystyle 1$ 个蓝球. 现有放回地从袋中取 $\displaystyle 3$ 次球, 每次取 $\displaystyle 1$ 个球, 则恰有 $\displaystyle 1$ 次取到蓝球 的概率为 \underline{\hspace{4em}}.}
% \ansat{198；【方法探究】 - 考点二：古典概型与几何概型 - 例 (1)}

% % --- 题目 ---
% \problem[P212-3 (07-1,3)]{设二维随机变量 $\displaystyle (X, Y)$ 的概率密度为
% \[ f(x, y) = \begin{cases} 2 - x - y, & 0 < x < 1, 0 < y < 1, \\ 0, & \text{其他}. \end{cases} \]
% \begin{enumerate}
% \item[(1)] 求 $\displaystyle P\{X > 2Y\}$;
% \item[(2)] 求 $\displaystyle Z = X + Y$ 的概率密度 $\displaystyle f_Z(z)$.
% \end{enumerate}}
% \ansat{361；【真题精选】 - 考点二：两个随机变量的函数的分布 - 3}

% % --- 题目 ---
% \problem[P221-5 (06-3)]{设随机变量 $\displaystyle X$ 的概率密度为
% \[ f_X(x) = \begin{cases} \frac{1}{2}, & -1 < x < 0, \\ \frac{1}{4}, & 0 < x < 2, \\ 0, & \text{其他.} \end{cases} \]
% 令 $\displaystyle Y = X^2$, $\displaystyle F(x, y)$ 为二维随机变量 $\displaystyle (X, Y)$ 的分布函数. 求:
% \begin{enumerate}
% \item[(1)] $\displaystyle Y$ 的概率密度 $\displaystyle f_Y(y)$;
% \item[(2)] $\displaystyle \mathrm{Cov}(X, Y)$;
% \item[(3)] $\displaystyle F\left(-\frac{1}{2}, 4\right)$.
% \end{enumerate}}
% \ansat{367；【真题精选】 - 考点二：随机变量的协方差与相关系数 - 5}

% % --- 题目 ---
% \problem[P219-9 (99-3)]{设随机变量 $\displaystyle X_{ij} \ (i, j = 1, 2, \dots, n; \ n \ge 2)$ 独立同分布, $\displaystyle EX_{ij} = 2$, 则行列式
% \[ Y = \begin{vmatrix} X_{11} & X_{12} & \dots & X_{1n} \\ X_{21} & X_{22} & \dots & X_{2n} \\ \vdots & \vdots & & \vdots \\ X_{n1} & X_{n2} & \dots & X_{nn} \end{vmatrix} \]
% 的数学期望 $\displaystyle EY = \underline{\hspace{4em}}.$}
% \ansat{366；【真题精选】 - 考点一：随机变量的数学期望与方差 - 9}

% --- 题目 ---
\problem[P204-3 (13-1)]{设随机变量 $\displaystyle X$ 的概率密度为 $\displaystyle f(x) = \begin{cases} \frac{1}{9}x^2, & 0 < x < 3, \\ 0, & \text{其他}. \end{cases}$, 
\newline
令随机变量 $\displaystyle Y = \begin{cases} 2, & X \le 1, \\ X, & 1 < X < 2, \\ 1, & X \ge 2. \end{cases}.$
\begin{enumerate}
\item[(1)] 求 \(Y\) 的分布函数;
\item[(2)] 求概率 \(P\left\{ X \leq Y \right\}\).
\end{enumerate}}
\ansat{357；【真题精选】 - 考点二：随机变量的函数的分布 - 3}

% % --- 题目 ---
% \problem[P210-1 (12-1)]{设随机变量 $\displaystyle X$ 与 $\displaystyle Y$ 相互独立, 且分别服从参数为 $\displaystyle 1$ 与参数为 $\displaystyle 4$ 的指数分布, 则 $\displaystyle P\{X < Y\} =~(~\quad~)$
% }
% \begin{tasks}(4)
%   \task $\displaystyle \frac{1}{5}.$
%   \task $\displaystyle \frac{1}{3}.$
%   \task $\displaystyle \frac{2}{5}.$
%   \task $\displaystyle \frac{4}{5}.$
% \end{tasks}
% \ansat{360；【真题精选】 - 考点一：二维随机变量的分布 - 1}

% % --- 题目 ---
% \problem[P214-7 (17-1)]{设随机变量 $\displaystyle X$ 的分布函数为 $\displaystyle F(x) = 0.5\Phi(x) + 0.5\Phi\left(\frac{x-4}{2}\right)$, 其中 $\displaystyle \Phi(x)$ 为标准正态分布函数, 则 $\displaystyle EX = \underline{\hspace{4em}}. $}
% \ansat{362；【十年真题】 - 考点一：随机变量的数学期望与方差 - 7}

% --- 题目 ---
\problem[P196-11 (18-3)]{设随机事件 $\displaystyle A, B, C$ 相互独立, 且}
\[ P(A) = P(B) = P(C) = \frac{1}{2}, \]
{则 $\displaystyle P(AC|A \cup B) = \underline{\hspace{4em}}.$}
\ansat{354；【十年真题】 - 考点一：概率的五大公式 - 11}
\vspace{4em}

% % --- 题目 ---
% \problem[P207-5 (16-1,3)]{设二维随机变量 $\displaystyle (X,Y)$ 在区域
% \[ D = \{(x,y) \ | \ 0 < x < 1, \ x^2 < y < \sqrt{x}\} \]
% 上服从均匀分布, 令
% \[ U = \begin{cases} 1, & X \le Y, \\ 0, & X > Y. \end{cases} \].
% \begin{enumerate}
% \item[(1)] 写出 $\displaystyle (X,Y)$ 的概率密度;
% \item[(2)] 问 $\displaystyle U$ 与 $\displaystyle X$ 是否相互独立? 并说明理由;
% \item[(3)] 求 $\displaystyle Z = U + X$ 的分布函数 $\displaystyle F(z)$.
% \end{enumerate}}
% \ansat{358；【十年真题】 - 考点二：两个随机变量的函数的分布 - 5}

% --- 题目 ---
\problem[P200-1 (23-3-局部)]{设随机变量 $\displaystyle X$ 的概率密度为 $\displaystyle f(x) = \frac{\mathrm{e}^x}{(1+\mathrm{e}^x)^2}, -\infty < x < +\infty$, 令 $\displaystyle Y = \mathrm{e}^X$.}
\begin{enumerate}
\item[(1)] 求 $\displaystyle X$ 的分布函数;
\item[(2)] 求 $\displaystyle Y$ 的概率密度.
\end{enumerate}
\begin{note}
  主要错的是 (2)
\end{note}
\ansat{355；【十年真题】 - 考点二：随机变量的函数的分布 - 1}

% % --- 题目 ---
% \problem[P223-变式1 (03-3)]{设总体 $\displaystyle X$ 服从参数为 2 的指数分布, $\displaystyle X_1, X_2, \dots, X_n$ 为来自总体 $\displaystyle X$ 的简单随机样本, 则当 $\displaystyle n \to \infty$ 时, $\displaystyle Y_n = \frac{1}{n} \sum_{i=1}^n X_i^2$ 依概率收敛于 \underline{\hspace{4em}}.}
% \ansat{368；【方法探究】 - 考点：大数定律、中心极限定理与切比雪夫不等式 - 变式1}
% \vspace{6em}

% % --- 题目 ---
% \problem[P224-4 (18-3)]{设 $\displaystyle X_1, X_2, \dots, X_n (n \ge 2)$ 为来自总体 $\displaystyle N(\mu, \sigma^2)$ ($\displaystyle \sigma > 0$) 的简单随机样本. 令 $\displaystyle \bar{X} = \frac{1}{n} \sum_{i=1}^n X_i, \ S = \sqrt{\frac{1}{n-1} \sum_{i=1}^n (X_i - \bar{X})^2}, S^* = \sqrt{\frac{1}{n} \sum_{i=1}^n (X_i - \mu)^2},$
% 则 ( \quad )
% \begin{tasks}(2)
%   \task $\displaystyle \frac{\sqrt{n}(\bar{X} - \mu)}{S} \sim t(n)$.
%   \task $\displaystyle \frac{\sqrt{n}(\bar{X} - \mu)}{S} \sim t(n-1)$.
%   \task $\displaystyle \frac{\sqrt{n}(\bar{X} - \mu)}{S^*} \sim t(n)$.
%   \task $\displaystyle \frac{\sqrt{n}(\bar{X} - \mu)}{S^*} \sim t(n-1)$.
% \end{tasks}}
% \ansat{369；【十年真题】 - 考点一：抽样分布 - 4}

% % --- 题目 ---
% \problem[P212-5 (01-3)]{设随机变量 $\displaystyle X$ 和 $\displaystyle Y$ 的联合分布是正方形 $\displaystyle G = \left\{ (x, y) | 1 \le x \le 3, 1 \le y \le 3 \right\}$ 上的均匀分布, 试求随机变量 $\displaystyle U = |X - Y|$ 的概率密度 $\displaystyle p(u)$.}
% \ansat{361；【真题精选】 - 考点二：两个随机变量的函数的分布 - 5}

% % --- 题目 ---
% \problem[P224-3 (23-1,3)]{设 $\displaystyle X_1, X_2, \dots, X_n$ 为来自总体 $\displaystyle N(\mu_1, \sigma^2)$ 的简单随机样本, $\displaystyle Y_1, Y_2, \dots, Y_m$ 为来自总体 $\displaystyle N(\mu_2, 2\sigma^2)$ 的简单随机样本, 且两样本相互独立. 记 $\displaystyle \bar{X} = \frac{1}{n} \sum_{i=1}^n X_i, \bar{Y} = \frac{1}{m} \sum_{i=1}^m Y_i$, $\displaystyle S_1^2 = \frac{1}{n-1} \sum_{i=1}^n (X_i - \bar{X})^2, \ S_2^2 = \frac{1}{m-1} \sum_{i=1}^m (Y_i - \bar{Y})^2,$ 则 ( \quad )
% \begin{tasks}(2)
%   \task $\displaystyle \frac{S_1^2}{S_2^2} \sim F(n, m)$.
%   \task $\displaystyle \frac{S_1^2}{S_2^2} \sim F(n-1, m-1)$.
%   \task $\displaystyle \frac{2S_1^2}{S_2^2} \sim F(n, m)$.
%   \task $\displaystyle \frac{2S_1^2}{S_2^2} \sim F(n-1, m-1)$.
% \end{tasks}}
% \ansat{369；【十年真题】 - 考点一：抽样分布 - 3}

% --- 题目 ---
\problem[P224-3 (24-3-仅数学三)]{设总体 $\displaystyle X$ 服从 $\displaystyle [0, \theta]$ 上的均匀分布, 其中 $\displaystyle \theta \in (0, +\infty)$ 为未知参数. $\displaystyle X_1, X_2, \dots, X_n$ 是来自总体 $\displaystyle X$ 的简单随机样本, 记 $\displaystyle X_{(n)} = \max\{X_1, X_2, \dots, X_n\}, T_c = cX_{(n)}$.
\begin{enumerate}
\item[(1)] 求 $\displaystyle c$, 使得 $\displaystyle E(T_c) = \theta$;
\item[(2)] 记 $\displaystyle h(c) = E[(T_c - \theta)^2]$, 求 $\displaystyle c$ 使得 $\displaystyle h(c)$ 最小.
\end{enumerate}}
\ansat{370；【十年真题】 - 考点二：统计量的数字特征 - 3}
